\chapter{}

Quando il Rosi era doventato padrone del Pesce Azzurro, c’era un ingresso solo, quello da Via dei Rossi, con un’insegna di ferro, a banderuola, ferma al muro e con un pesce dipinto tanto dall’una parte che dall’altra. Sulla porta, una Madonna in bassorilievo; del quattrocento. Ci stava ancora il lume attaccato, ma la fune per tirarlo giù mancava.

Poi furono aperti anche due ingressi dalla Via Cavour. E ad uno di questi, dietro il cristallo della porta, una vetrina a due piani, foderata con la carta che cambiavano una volta tutte le settimane; piena di polli già pelati, di carni arrostite, e d’altre delizie.

Dopo l’ingresso da Via dei Rossi una gran porta, per entrare in una piazzola interna sempre ingombra di calessi e d’ogni specie di legni. Accanto a questa, la stalla; che poteva contenere fino a trenta bestie. Sopra la stalla, la capanna.

Tutti i sabati, Domenico faceva l’elemosina dei pezzi di pane avanzati agli avventori.

La stretta Via dei Rossi, al principio, dov’era l’uscio vecchio della trattoria, si empiva un’ora prima del tempo, di mendicanti; fra i quali era anche la moglie di Pipi, giovine, ma così smunta e gialla che la sua bocca era come un taglio senza labbra: andava come se non avesse potuto piegare la testa da nessuna parte. Molte volte, dalla veste male abbottonata e sudicia, si vedeva il petto vuoto e senza i seni.

C’era anche una vecchia, dal naso enorme e pavonazzo, con un cappello da contadina, del quale le trecce di paglia si disfacevano intorno; e ne rimaneva sempre un giro di meno. Questa pretendeva d’avere la prima elemosina, e non se ne andava finchè tutti i pezzi di pane non fossero stati distribuiti. Talvolta gridava:

— Quella vecchiaccia ne ha avuto più di me.

Ed apriva ancora i lembi del fazzoletto pieno di pane duro, sorreggendo sotto l’ascella il bastoncino.

C’era una mendicante, a cui Domenico faceva l’elemosina tre giorni della settimana; una donna grande, dal volto acceso ed uguale come una maschera sottile, che non si poteva togliere, una maschera di pelle rossa. Portava, d’estate e d’inverno, uno scialletto di lana nera annodato dietro il dorso. Teneva sempre incrociate le mani pallide sul petto. La sua figliola, alta e leggiadra, non la lasciava mai, tenendo una mano infilata sotto uno dei suoi bracci; era scema e sorrideva sempre; ma di un sorriso dolce ed appassionato.

Camminavano ambedue rasentando i muri; a passi lunghi, come se avessero voluto fuggire. Nell’attraversare la strada da una parte all’altra, si affrettavano anche di più.

Quando mangiavano la zuppa a qualche convento, la figliola voltava il dorso a tutti; e ritraendo il cucchiaio dalla bocca, faceva grandi risate silenziose.

Quando la madre morì, fu rinchiusa in un manicomio.

C’era un cieco, che imprecava contro il figlio; che aveva una mano secca con un dito di meno:

— Sei un mascalzone, e non mi aiuti. Se tu stai costi appoggiato al muro, non troverai più pane per noi. Mascalzone! Mascalzone!

E tendeva un orecchio, accartocciandovi dietro una mano; per capire quanta elemosina ci fosse ancora; mentre la voce era la stessa di quando recitava le devozioni.

Tutti gli altri poveri erano andati incontro a Rosaura come un branco di polli verso il punto dov’è rimbalzato un chicco di granturco.

Il giovinetto del cieco ascoltava, scalcinando con le dita le commessure dei mattoni: preferiva esser l’ultimo perchè, senza leticare, era sicuro che Rosaura avrebbe serbato qualche cosa per lui.

Tutte le mendicanti guardavano il pane avuto; e qualcuna ne riposava un pezzo troppo secco dentro una fenditura del muro, che era accanto all’uscio. Allora Rosaura, sporgendosi tutta fuori, esclamava:

— Guardatela: viene a chiedere l’elemosina, e poi la scrafia!

Una donna rispondeva, tenendosi ambedue le mani strette sopra i fianchi:

— Se l’avessi avuta io, l’avrei mangiato!

Qualcuna rideva, addentando il pane; dopo averlo un poco rigirato tra le mani sudicie. Ad un tratto, dal mormorio basso e incomprensibile, cominciava un alterco:

— Viene a chiedere il pane, ed è ricca quanto vuole.

— Che importa a te? Sono ricca?... Non le dia retta.

Rosaura interrompeva:

— State zitta, altrimenti non ve ne daremo più.

Un’altra donna, con il volto guasto da un ezzema, bendato con una pezzuola azzurra annodata dietro la testa, rispondeva:

— Ha ragione. Ma io non mi sono mai lamentata.

Si vedevano soltanto i suoi occhi infiammati, come piaghe, che non potevano stare aperti; ed era costretta, per guardare, a sollevare il capo di traverso; mentre, parlando, la benda seguiva i movimenti della bocca. E che bocca aveva!

Un vecchio, che sopravveniva quasi sempre a elemosina finita, cercava d’impietosire con quel tono che i mendicanti adoprano:

— Per amor di Dio.... anche a me.

— Non c’è più niente. Perchè non venite prima?

— Le gambe non mi reggono più!

E batteva il suo bastone su lo scalino dell’uscio. Rosaura se ne andava senza dargli niente; dopo avergli risposto:

— Ma per arrivare ora vi reggono!

Allora egli aspettava ancora per lungo tempo; con un’ostinazione rabbiosa:

— Signora mia, non mi faccia soffrire più!

Aveva lavorato tutta la vita; e pensava, come a una magnificenza, che se si ammalasse avrebbe potuto entrare in un ospedale, dove sarebbe stato tutto il giorno steso sopra il letto. E a mangiare bene!

La moglie almeno gli era morta giovine, e non soffriva più! Ma egli finì con il credere un obbligo l’elemosina, come trovare uno scalone e mettercisi a sedere senza che lo mandassero via.
